\section{Introduction}
\label{sec:intro}

Large language models sometimes lie: they state what they believe to be false
\citep{park2023deception}, to conceal misconduct \citep{scheurer2023strategic,jarviniemi2024uncovering} or
to fake alignment \citep{greenblatt2024alignment}. Following Liars' Bench \citep{kretschmar2025liars},
we call a statement a lie when the model asserts something it believes to be false, belief being measured
by what it consistently asserts under no pressure to lie; this separates lying from hallucination. The
cheapest detector of such lies is a linear probe on the residual stream
\citep{alain2016understanding,azaria2023internal,marks2023geometry}: one dot product per response, no
extra forward pass, and a chance to catch a lie that leaves no trace in the transcript
\citep{goldowskydill2025detecting}. Its weakness is generalisation. A probe trained on one distribution of
lies often fails on others \citep{kirch2025offpolicy,kumar2026pressure,natarajan2026targeted}, and Liars'
Bench makes the failure concrete: across seven datasets of on-policy lies from four open-weight models,
probes trained on instruction pairs are near chance in balanced accuracy at the benchmark's threshold on
most of them, and on the hardest, the harm-pressure knowledge report (\hpkr), in which a model denies
knowing the answer to a hazardous question it has been shown to know, the benchmark's probe scores an
AUROC of \V{lb-mean-auroc-hpkr} on Llama-3.3-70B-Instruct at its layer~20, worse than chance. At the
model's middle layer a probe trained the same way ranks \hpkr{} well, yet still flags none of its lies.

\citet{mahadik2026persona} propose persona coordinates as a remedy: the model role-plays deceptive and
honest personas, each persona's mean activation is a vector, and the difference between the two group
means is a direction that uses no evaluation data and so cannot fit one dataset's idiosyncrasies. On
Llama~3.2-3B and Llama~3-8B these coordinates transfer better than probes on raw activations; the
authors name larger models and richer persona sets as future work. We take up both, and find that the
published contrast, rebuilt for Llama-3.3-70B-Instruct, inverts on \hpkr{} (AUROC \V{zs-dec-hpkr-L20} at
the benchmark's layer). Every \hpkr{} lie is a denial of knowledge, while every published deceptive
persona lies by controlling disclosure. Adding personas that deny knowledge they have raises the
contrast's ranking of \hpkr, but for the wrong reason: the new direction flags nearly every reply that says ``I do not know'', true or false, and where a lie
and an honest answer are the same sentence every persona direction is within reach of random directions.
A persona direction reads the register of a lie, not what makes it a lie.

The definition of a lie offers a second route. A lie about knowledge is a denial by a model that knows,
and what a model knows can be measured without any labelled lie, from questions it knows and questions
nobody can know. We build a \emph{knowledge direction} from exactly that, with the same denial forced
after each question so that the reply's words do not vary, and ask whether a detector made of a persona
contrast and a knowledge direction transfers across Liars' Bench better than supervised probes, whether
it separates lies from honest answers that use the same words, what each part reads, whether the result
survives resampling, new questions, other layers and strict thresholds, and whether it holds on other
models and benchmarks.

The answer is mixed, and we state it up front. The persona contrast is the general half: at each model's
middle layer it transfers better than supervised probes on three of the benchmark's four models, and on
the main model it beats random contrasts of a random bank, but it fails on Qwen, even with a bank Qwen
plays far more fully. The knowledge direction holds where a lie is a denial in an honest answer's words,
on two models and narrowly a third; beyond those slices its transfer rests on the benchmark's request
wording, and its detection of a fine-tuned liar holds on one model and not on the other. Their sum ranks
the benchmark's lies above a probe trained at the same depth on all four models (on Qwen both are near
chance) but is not significantly better than the persona contrast on the main model's average, and on an
external competition it sits mid-table behind a supervised probe. Our contributions are:
\begin{itemize}[leftmargin=1.2em, itemsep=0pt, topsep=1pt]
\item a persona bank for deception by denial and a single label-free contrast from it that transfers
  better than supervised probes on three of four models, checked against an unrelated widening and a
  random bank, whose test of its gain over the published contrast it does not pass
  (\cref{sec:persona,sec:models,sec:external});
\item a register diagnosis: on lies about knowledge, persona directions read the words of a denial
  (\cref{sec:register});
\item a knowledge direction with its boundaries measured: same-word separation on two models (narrowly a
  third), a small margin on the one fine-tune the text cannot separate, and only the wording-controlled
  results free of the benchmark's wording (\cref{sec:knowledge,sec:models});
\item a combined detector, above a same-depth probe on all four models, level with the persona contrast
  on the main model, mixed across models and weaker in the tail (\cref{sec:combined}), with a causal
  check of its persona part by steering (\cref{sec:steering});
\item an evaluation protocol for zero-shot directions (random bands, a random bank, wording-controlled
  slices, the base model's read as a control for organisms, paired intervals) and two pre-registered
  external tests reported with their failures (\cref{sec:external}).
\end{itemize}
